\chapter*{Abstract}
\addcontentsline{toc}{chapter}{Abstract}

Large-scale manufacturing organisations such as IKEA of Sweden rely on standardised
product label templates to communicate product identity, regulatory compliance, and
logistical information across global supply chains. Ensuring that every generated label
exactly conforms to its approved baseline design is currently performed by human quality
assurance testers—a process that is slow, costly, and error-prone at production volume.

This thesis presents the design, implementation, and evaluation of an automated hybrid
validation system for standardised label templates. The system accepts a structured
JSON metadata payload and a NiceLabel (\texttt{.nlbl}) label file as inputs, renders the
generated label into a pixel-accurate image, and applies a three-tier validation pipeline:
(1) rule-based validation of structured metadata fields, (2) reference-based visual
comparison against approved baseline templates using image similarity metrics, and (3) a
machine-learning-based binary classifier that captures structural deviation patterns
beyond simple pixel differences.

Three visual comparison methods were implemented and evaluated on a dataset of 150
labelled template images spanning 20 template types, with controlled synthetic violations:
a Structural Similarity Index Measure (SSIM) baseline, an ORB feature-matching approach,
and a Residual Network (ResNet-50) binary classifier. A Siamese neural network was also
trained to learn pairwise similarity directly from template image pairs. The final hybrid
system, combining SSIM thresholding, ORB spatial matching, and the ResNet-50 classifier
in a majority-vote ensemble, achieved a validation accuracy of 93.3\%, a precision of
92.8\%, a recall of 91.2\%, and an F1-score of 92.0\% on the held-out test set, with a
median processing time of 347\,ms per label. Violation localisation using bounding-box
overlay achieved an Intersection over Union (IoU) above 0.5 in 78.1\% of detected errors.

The results demonstrate that a hybrid approach substantially outperforms any single
comparison method and that the Siamese and ResNet classifiers capture structural
violations that pixel-level metrics miss. The system produces an actionable, explainable
report that identifies \emph{which} layout elements violate the reference and \emph{where}
on the label the violation occurs, directly supporting manual correction workflows.

\vspace{0.8em}
\noindent\textbf{Keywords:}
document layout analysis, automated template validation, structural similarity,
convolutional neural network, industrial quality control, visual document understanding,
Siamese network, object detection, reference-based comparison, NiceLabel.
\clearpage
